%laden der Präambel mit Latexbefehlen/-klassen
\include{bericht.preambel}

% DATEN AKTUALISIEREN!
% Ort der Exkursion
\def\VARort{HKW Reuter West}
% Datum der Exkursion
\def\VARdatum{27. Januar 2022}
% Autor*in: Vorname und Name 
\def\VARname{Jane Doe}
% Autor*in: Studiengang
\def\VARstudiengang{Rocket Science}
% Autor*in: Matrikel-Nr.
\def\VARmatrikelnr{111312}

\begin{document}
%globale Einstellung für Quellcodedarstellung mit listings
\lstset{basicstyle=\scriptsize\ttfamily,language={[LaTeX]TeX}}

%Seitennummerierung arabisch
\pagenumbering{arabic}

%Titel
\input{00a_titelei}
%Abstract
\input{00b_abstract}
%Inhaltsverzeichnis
\setcounter{tocdepth}{2}
\tableofcontents

%Beginn der Kapitel
\input{01_einleitung}
	
\input{02_latex}

%weitere Kapitel hier einfügen
%\input{03_neues_kapitel}

%Literatur
\clearpage
\printbibliography[heading=bibintoc]

%Bildnachweise
%Kann auskommentiert werden, wenn keine Übernahme von Abbildungen erfolgt ist.
\clearpage
\input{97_bildnachweis}

%%Abkürzungs- und Symbolverzeichnis
\clearpage
\input{98_symbole_abk}

%Abbildungsverzeichnis
\clearpage
\addcontentsline{toc}{section}{\listfigurename}
\listoffigures

%Tabellenverzeichnis
\clearpage
\addcontentsline{toc}{section}{\listtablename}
\listoftables

%Anhang
%Ausführliche Datentabellen, umfangreiche Fließbilder und lange Berechnungen können in den Anhang
%\clearpage
%\appendix
%\input{99_anhang}

\end{document}